\clearpage
\VerdierPhysicalPageAnchor{0231}%
\begingroup\small\setlength{\parskip}{0pt}\setlength{\emergencystretch}{1.5em}
\setlength{\abovedisplayskip}{4pt plus 2pt minus 2pt}%
\setlength{\belowdisplayskip}{4pt plus 2pt minus 2pt}%

\VerdierSemanticAnchor{ch4d-4-1-7-proof}
Supposons, tout d'abord, la proposition démontrée lorsque $X$ est un double
complexe et lorsque l'entier $m=1$, et montrons comment on peut se ramener
à ce cas. Soit $\varphi\colon[n]\longrightarrow[2]$ l'application :

\VerdierSemanticAnchor{ch4d-4-1-7-1}
\begin{equation}
\left\{
\begin{aligned}
\varphi(m)&=1,\\
\varphi(i)&=2,\qquad i\in[n],\quad i\neq m.
\end{aligned}
\right.
\tag{4.1.7.1}
\end{equation}

Posons $\int_\varphi X=Z$ (chap.~I, 2.2.2). On peut alors effectuer sur le
double complexe $Z$ les constructions décrites en (4.1.6.$i$),
$1\leq i\leq4$. On obtient des doubles complexes :

\VerdierSemanticAnchor{ch4d-4-1-7-2}
\begin{equation}
{}^{1}\!Z(p,q),\qquad
(p,q)\in\mathcal{F}\ell_1(\widetilde{\mathbb Z}),
\tag{4.1.7.2}
\end{equation}

et des morphismes de doubles complexes :

\VerdierSemanticAnchor{ch4d-4-1-7-3}
\begin{equation}
\begin{aligned}
f_Z\bigl((p,q),(p',q')\bigr)&\colon
{}^{1}\!Z(p,q)\longrightarrow{}^{1}\!Z(p',q'),\\
g_Z(p,q,r)&\colon
{}^{1}\!Z(q,r)\longrightarrow T_1\bigl({}^{1}\!Z(p,q)\bigr).
\end{aligned}
\tag{4.1.7.3}
\end{equation}

On constate alors que :

\VerdierSemanticAnchor{ch4d-4-1-7-4}
\begin{equation}
\left\{
\begin{aligned}
\int_\varphi{}^{m}\!X(p,q)
  &={}^{1}\!Z(p,q),
  &&(p,q)\in\mathcal{F}\ell_1(\widetilde{\mathbb Z}),\\
\int_\varphi f\bigl((p,q),(p',q')\bigr)
  &=f_Z\bigl((p,q),(p',q')\bigr),
\end{aligned}
\right.
\tag{4.1.7.4}
\end{equation}

(vérification immédiate) et que de plus le morphisme $g_Z(p,q,r)$ est le
morphisme composé :

\VerdierSemanticAnchor{ch4d-4-1-7-5}
\begin{equation}
\int_\varphi{}^{m}\!X(q,r)
\xrightarrow{\int_\varphi g(p,q,r)}
\int_\varphi T_m\bigl({}^{m}\!X(p,q)\bigr)
\xrightarrow{(**)}
T_1\left(\int_\varphi{}^{m}\!X(p,q)\right)
\tag{4.1.7.5}
\end{equation}

où $(**)$ est l'isomorphisme de commutation du foncteur $\int_\varphi$ aux
translations (chap.~I, 2.3.2.1) (pour vérifier ceci, on utilise la définition
explicite des différentielles de $Z$ (chap.~I, 2.2.2.5) et de l'isomorphisme
de commutation (chap.~I, 2.3.2.1). Construisons alors l'objet spectral défini
par $Z$ filtré par son premier degré. Il résulte de l'égalité
$\int_2\int_\varphi=\int_n$, de (4.1.7.4), (4.1.7.5) et de
(chap.~I, 2.3.3) que
$\bigl({}^{m}\!Y(p,q),{}^{m}\!d(p,q,r)\bigr)$ est égal à cet objet
spectral.
\endgroup
